\chapter{Минимальная конечная алгебра и происхождение кратности}

Рассмотрим минимальную конечную алгебру с двумя примитивными центральными
идемпотентами
\begin{equation}
 A_F=\mathbb C\oplus\mathbb C,
 \qquad
 p_X=(1,0),
 \qquad
 p_C=(0,1).
\end{equation}
Её минимальное точное представление на \(H_F=\mathbb C^2\) имеет
\begin{equation}
 \rho(a_X,a_C)=
 \begin{pmatrix}a_X&0\\0&a_C\end{pmatrix},
 \qquad
 \Tr p_X=\Tr p_C=1.
\end{equation}
Таким образом, кратность два и равные конечные trace-веса возникают из
минимального faithful representation, а не из таблицы секторных
коэффициентов.

\section{Канонический факторный модуль}

Пусть \(\widehat1_X=\pi^{-1}\) и
\(\widehat1_C=(2\pi)^{-1/2}\) --- нормированные постоянные моды факторов.
Внутри \(L^2(K)\otimes H_F\) определим спектрально выделенный модуль
\begin{equation}
 \mathcal H_{\mathrm{fac}}
 =
 \bigl(L^2(X)\otimes\mathbb C\widehat1_C\otimes p_XH_F\bigr)
 \oplus
 \bigl(\mathbb C\widehat1_X\otimes L^2(C)\otimes p_CH_F\bigr).
\end{equation}
Проекторы на \(\widehat1_X,\widehat1_C\) являются каноническими
нулевыми спектральными проекторами лапласианов связных факторов.

Для единичных факторных мод полный \(L^2(K)\)-след даёт
\begin{align}
 \|1_X\otimes\widehat1_C\otimes p_X\|^2
 &=\pi^2,\\
 \|\widehat1_X\otimes1_C\otimes p_C\|^2
 &=2\pi.
\end{align}
Следовательно,
\begin{equation}
 \|\Xi_{\mathrm{fac}}\|^2=\pi^2+2\pi.
\end{equation}

\begin{theorem}[Minimal finite-module pass]
Алгебра \(\mathbb C\oplus\mathbb C\), её минимальное faithful
представление и нулевые spectator-проекторы канонически выводят две
факторные компоненты, их единичные trace-веса и требуемые собственные
объёмы. Ни независимые sector weights, ни выбранные точки дефектов не
требуются.
\end{theorem}

\section{Локальность и спектральный статус}

Редуцированный оператор
\begin{equation}
 D_{\mathrm{fac}}=D_X\oplus D_C
\end{equation}
на \(\mathcal H_{\mathrm{fac}}\) имеет компактную резольвенту, поскольку оба
фактора компактны. Однако \(\mathcal H_{\mathrm{fac}}\) является
спектрально спроецированным подпространством полного \(L^2(K)\), а не
пространством всех локальных скалярных полей на \(K\). Поэтому конструкция
является допустимым spectral reduction, но не обычной локальной EFT на
четырёхмерном произведении.

\section{First-order red-team минимального смешивания}

Проверим простейший конечный оператор
\begin{equation}
 D_F=\begin{pmatrix}0&m\\\overline m&0\end{pmatrix},
 \qquad
 J_F=\text{комплексное сопряжение}.
\end{equation}
Для \(a=(a_X,a_C)\), \(b=(b_X,b_C)\) условие первого порядка требует
\begin{equation}
 [[D_F,\rho(a)],J_F\rho(b)J_F^{-1}]=0.
\end{equation}
Его внедиагональная компонента пропорциональна
\begin{equation}
 m(a_C-a_X)(b_C-b_X).
\end{equation}
Для произвольных \(a,b\in A_F\) это условие вынуждает \(m=0\).

\begin{nogocriterion}[Minimal diagonal first-order no-go]
В минимальном диагональном представлении
\(\mathbb C\oplus\mathbb C\) простая real structure и условие первого
порядка запрещают ненулевое смешивание двух листов. Модель выводит
кинематический модуль и норму, но не выводит vertex, связывающую обе
факторные компоненты в одну charged-lepton динамику.
\end{nogocriterion}

\section{Текущий статус parent-action программы}

Впервые закрыты одновременно:
\begin{enumerate}
 \item непрерывный масштаб факторных фаз;
 \item относительные факторные меры;
 \item происхождение кратности два;
 \item равенство конечных trace-весов.
\end{enumerate}
Открытым остаётся не нормировочный, а операторный вопрос: построить
реальный градуированный бимодуль, в котором ненулевой odd operator
удовлетворяет first-order condition и порождает физическую вершину без
новых непрерывных коэффициентов. Следующий гейт должен классифицировать
минимальные \(A_F\)-бимодули, а не менять геометрическую меру.

Машинный аудит сохранён в
\texttt{s2t\_v3\_minimal\_finite\_algebra\_results.json}.